%! TeX program = lualatex
\documentclass[../main.tex]{subfiles}

\begin{document}

\section{Week 1 practice problems}

\begin{exercise}
  For each log-log plot below, determine its power rule model in \(y = c x^{\alpha}\) and in \(y \propto x^{\alpha}\) forms.

  \begin{tikzpicture}
    \begin{axis}[loglog]
      \addplot[very thick] {3*x + 1};
    \end{axis}
  \end{tikzpicture}
  \hfill
  \begin{tikzpicture}
    \begin{axis}[loglog]
      \addplot[very thick] {-1/2*x + 2};
    \end{axis}
  \end{tikzpicture}

  \begin{tikzpicture}
    \begin{axis}[loglog]
      \addplot[very thick] {3/2*x + 3};
    \end{axis}
  \end{tikzpicture}
  \hfill
  \begin{tikzpicture}
    \begin{axis}[loglog]
      \addplot[very thick] {-1/4*x + 2};
    \end{axis}
  \end{tikzpicture}

  \begin{solution}
    \phantom{top}

    \begin{itemize}
      \item Top left.  \(y = e x^{3}\) and \(y \propto x^{3}\).

      \item Top right. \(y = e^{2} x^{-1/2}\) and \(y \propto x^{-1/2}\).

      \item Bottom left. \(y = e^{3} x^{3/2}\) and \(y \propto x^{3/2}\).

      \item Bottom right. \(y = e^{2} x^{-1/4}\) and \(y \propto x^{-1/4}\).
    \end{itemize}
  \end{solution}
\end{exercise}

\begin{exercise}
  Graph the log-log transformation of the power rule models. 
  \begin{multicols}{4}
    \begin{itemize}
      \item \(y = x^{1}\)
      \item \(y = x^{2}\)
      \item \(y = x^{3}\)
      \item \(y = \frac{1}{2}x^{1}\)
      \item \(y = \frac{1}{2}x^{2}\)
      \item \(y = \frac{1}{2}x^{3}\)
      \item \(y = x^{-1}\)
      \item \(y = x^{-2}\)
      \item \(y = x^{-3}\)
      \item \(y = \frac{1}{2}x^{-1}\)
      \item \(y = \frac{1}{2}x^{-2}\)
      \item \(y = \frac{1}{2}x^{-3}\)
    \end{itemize}
  \end{multicols}

  For each model, reflect on the following questions.
  \begin{enumerate}
    \item Does \(x\) and \(y\) exhibit isometry or allometry?
    \item If \(x\) increases, does \(y\) increase or decrease? Which feature of the plot shows this trend?
    \item If \(x\) decreases, does \(y\) increase or decrease? Which feature of the plot shows this trend?
    \item If \(y\) increases, does \(x\) increase or decrease? Which feature of the plot shows this trend?
    \item If \(y\) decreases, does \(x\) increase or decrease? Which feature of the plot shows this trend?
  \end{enumerate}

  \begin{solution}
    \phantom{top}

    \begin{tikzpicture}
      \begin{axis}[loglog line, title={\(y = x^{1} \rightarrow \ln(y) = \ln(x)\)}]
        \addplot[very thick] {x};
      \end{axis}
    \end{tikzpicture}
    \begin{tikzpicture}
      \begin{axis}[loglog line, title={\(y = \frac{1}{2}x^{1} \rightarrow \ln(y) = \ln(x) - \ln(2)\)}]
        \addplot[very thick] {x - ln(2)};
      \end{axis}
    \end{tikzpicture}


    \begin{tikzpicture}
      \begin{axis}[loglog line, title={\(y = x^{2} \rightarrow \ln(y) = 2\ln(x)\)}]
        \addplot[very thick] {2*x};
      \end{axis}
    \end{tikzpicture}
    \begin{tikzpicture}
      \begin{axis}[loglog line, title={\(y = \frac{1}{2}x^{2} \rightarrow \ln(y) = 2\ln(x) - \ln(2)\)}]
        \addplot[very thick] {2*x - ln(2)};
      \end{axis}
    \end{tikzpicture}

    \begin{tikzpicture}
      \begin{axis}[loglog line, title={\(y = x^{3} \rightarrow \ln(y) = 3\ln(x)\)}]
        \addplot[very thick] {3*x};
      \end{axis}
    \end{tikzpicture}
    \begin{tikzpicture}
      \begin{axis}[loglog line, title={\(y = \frac{1}{2}x^{3} \rightarrow \ln(y) = 3\ln(x) - ln(2)\)}]
        \addplot[very thick] {3*x - ln(2)};
      \end{axis}
    \end{tikzpicture}


    \begin{tikzpicture}
      \begin{axis}[loglog line, title={\(y = x^{-1} \rightarrow \ln(y) = -\ln(x)\)}]
        \addplot[very thick] {-x};
      \end{axis}
    \end{tikzpicture}
    \begin{tikzpicture}
      \begin{axis}[loglog line, title={\(y = \frac{1}{2}x^{-1} \rightarrow \ln(y) = -\ln(x) - \ln(2)\)}]
        \addplot[very thick] {-x - ln(2)};
      \end{axis}
    \end{tikzpicture}

    \begin{tikzpicture}
      \begin{axis}[loglog line, title={\(y = x^{-2} \rightarrow \ln(y) = -2\ln(x)\)}]
        \addplot[very thick] {-2*x};
      \end{axis}
    \end{tikzpicture}
    \begin{tikzpicture}
      \begin{axis}[loglog line, title={\(y = \frac{1}{2}x^{-2} \rightarrow \ln(y) = -2\ln(x) - \ln(2)\)}]
        \addplot[very thick] {-2*x - ln(2)};
      \end{axis}
    \end{tikzpicture}

    \begin{tikzpicture}
      \begin{axis}[loglog line, title={\(y = x^{-3} \rightarrow \ln(y) = -3\ln(x)\)}]
        \addplot[very thick] {-3*x};
      \end{axis}
    \end{tikzpicture}
    \begin{tikzpicture}
      \begin{axis}[loglog line, title={\(y = \frac{1}{2}x^{-3} \rightarrow \ln(y) = -3\ln(x) - ln(2)\)}]
        \addplot[very thick] {-3*x - ln(2)};
      \end{axis}
    \end{tikzpicture}
  \end{solution}
\end{exercise}



\begin{exercise}
  We explore allometric relations with biological models. 

  \begin{figure}[H]
    \centering

    \begin{tabular}{l|l|l}
      Quantity & Quantity & Scaling Relation \\ \midrule\midrule
      Metabolic rate (\(R\)) & Body mass ($m$) & $R \propto m^{3/4}$ \\ \midrule
      Heart rate (\(H\)) & Body mass ($m$) & $H \propto m^{-1/4}$
    \end{tabular}
  \end{figure}

  \begin{enumerate}
    \item If \sout{heart rate} mass increases by \(5\%\), how does metabolic rate change in terms of scaling and percentage?
      \begin{solution}
        The model is \(R \propto m^{3/4}\). A \(5\%\) increase means we scale mass by \(1.05\). In terms of scaling, metabolic rate scales \emph{by a factor} of \(1.05^{3/4} \approx 1.037\) which is an increase. In terms of percentage, metabolic rate increase by roughly \(3.7\%\).
      \end{solution}

    \item If \sout{heart rate} mass decreases by \(7\%\), how does metabolic rate change in terms of scaling and percentage?
      \begin{solution}
        The model is \(R \propto m^{3/4}\). A \(7\%\) decrease means we scale mass by \(0.93\). In terms of scaling, metabolic rate scales \emph{by a factor} of \(0.93^{3/4} \approx 0.947\) which is an decrease. In terms of percentage, metabolic rate decreases by roughly \(1 - 0.947 = (5.3\%)\).
      \end{solution}

    \item If body mass increases by \(10\%\), how does heart rate change in terms of scaling and percentage?
      \begin{solution}
        The model is \(R \propto m^{-1/4}\). A \(10\%\) increase means we scale mass by \(1.1\). In terms of scaling, heart rate scales \emph{by a factor} of \(1.1^{-1/4} \approx 0.976\) which is an decrease. In terms of percentage, heart rate decreases by roughly \(1 - 0.976 = (2.4\%)\).

        \hlwarn{Correction: the model should read \(H \propto m^{-1/4}\) (heart rate, not \(R\)).}
      \end{solution}

    \item If body mass decreases by \(5\%\), how does heart rate change in terms of scaling and percentage?
      \begin{solution}
        The model is \(R \propto m^{-1/4}\). A \(5\%\) decrease means we scale mass by \(0.95\). In terms of scaling, heart rate scales \emph{by a factor} of \(0.95^{-1/4} \approx 1.0129\) which is an increase. In terms of percentage, heart rate increases by roughly \(1.29\%\).

        \hlwarn{Correction: the model should read \(H \propto m^{-1/4}\) (heart rate, not \(R\)).}
      \end{solution}

    \item If a patient's heart rate increases, do the models predict their metabolic rate to increase or decrease?
      \begin{solution}
        We want heart rate to be the independent variable, so we isolate for \(H\) in terms of \(m\) to get \(m \propto R^{-1/4}\). Plug \(m\) into the metabolic rate vs body mass model to get \(R \propto (H^{-4})^{3/4}\). Simplify the exponents to get \(R \propto H^{-3}\).

        The exponent on heart rate is negative. Therefore, the model predicts that an increase in heart rate results in a decrease in metabolic rate.

        \hlwarn{Correction: we isolate for \(m\) in terms of \(H\) to get \(m \propto H^{-4}\) (not \(m \propto R^{-1/4}\)); the substitution then correctly gives \(R \propto (H^{-4})^{3/4} = H^{-3}\).}

      \end{solution}

    \item If a patient's metabolic rate decreases, do the models predict their heart rate to increase or decrease?

      \begin{solution}
        We want metabolic rate to be the independent variable, so we isolate for \(R\) in terms of \(m\) to get \(m \propto R^{4/3}\). Plug \(m\) into the heart rate vs body mass model to get \(H \propto (R^{4/3})^{-1/4}\). Simplify the exponents to get \(H \propto R^{-1/3}\).

        The exponent on metabolic rate is negative. Therefore, the model predicts that an decrease in metabolic rate results in an increase in heart rate.

        \hlwarn{Correction: the phrasing should be ``isolate for \(m\) in terms of \(R\)''; the resulting equation \(m \propto R^{4/3}\) is correct.}
      \end{solution}
  \end{enumerate}
\end{exercise}

\begin{exercise}
  Consider scaling relation in drug dosage.
  \begin{enumerate}
    \item Suppose the dosage of a particular drug and patient's body weight scale isometrically.  If a patient weighs \(150\) lbs should receive \(100\)mg of dosage, what dosage should a \(170\) lbs patient receive?

      \begin{solution}
        We let \(m\) denote body weight (lbs) and \(d\) drug dosage (mg).  Being isometrically related means \(d = (\text{constant}) \times m\). Given \(m = 150\) lbs and \(d = 100\) mg, we can solve for \[\text{constant} = \frac{d}{m} = \frac{100}{150}.\]  Therefore, the model is \(d = \frac{100}{150} m\), and an \(170\) lbs patient should receive a dosage of \[d = \frac{100}{150} 170 \approx 113 \text{ mg}.\]
      \end{solution}

    \item Suppose the dosage of a particular drug and patient's body weight scale allometrically.  A patient has been receiving \(100\)mg of dosage. If the patient experiences a \(5\%\) weight loss, and a new prescription says the patient should receive \(95\) mg (of the same) drug, should you suspect there is a mistake?

      \begin{solution}
        Yes, there could be a mistake. The proposed drug dosage in relation to weight change \emph{coincides} with an isometric relation, which contradicts the given data. 
      \end{solution}
  \end{enumerate}
\end{exercise}


\begin{exercise}
  Electrical power (\(P\)), voltage (\(V\)) and resistance (\(R\)) are related by \(\ln(P) = -\ln(R) + 2\ln(V)\). 
  \begin{enumerate}
    \item If \(V\) remains constant and \(R\) increases by \(100\%\), how does \(P\) change?

      \begin{solution}
        Convert the log-log transformation back into \(P = \frac{V^{2}}{R}\).

      If \(V\) remains constant, then \(P \propto R^{-1}\). An \(100\%\) increase in \(V\) means \(V\) scales by \(2\). Therefore \(P\) scales by a factor of \(2^{-1} = 1/2\) which is a decrease.  In terms of percentage, \(P\) decreases by \(50\%\).

      \hlwarn{Correction: the \(100\%\) increase is in \(R\), so \(R\) scales by \(2\); with \(P \propto R^{-1}\), \(P\) then scales by \(2^{-1}=1/2\).}
      \end{solution}

    \item If \(R\) remains constant and \(V\) doubles, how does \(P\) change?

      \begin{solution}
        Convert the log-log transformation back into \(P = \frac{V^{2}}{R}\).

        If \(R\) remains constant, then \(P \propto V^{2}\). Doubling \(V\) means \(V\) scales by \(2\). Therefore, \(P\) scales by \(2^{2} = 4\) which is an increase.  In terms of percentage, \(P\) increases by \(300\%\).
      \end{solution}

      
    \item If \(V\) remains constant and \(P\) halves, how does \(R\) change?
      \begin{solution}
        Convert the log-log transformation back into \(P = \frac{V^{2}}{R}\).

        The question wants information on \(R\), so we want \(R\) to be the dependent variable.  The equation rewrites as \(R = \frac{V^{2}}{P}\). Since \(V\) remains constant in this scenario, then \(R \propto P^{-1}\). Having \(P\) means \(P\) scales by \(0.5\). Therefore, \(R\) scales by \(0.5^{-1} = 2\) which is an increase.  In terms of percentage, \(R\) increases by \(100\%\).

        \hlwarn{Correction: ``Having \(P\)'' should read ``Halving \(P\)'' --- \(P\) scales by \(0.5\).}
      \end{solution}
  \end{enumerate}
\end{exercise}
\end{document}
